\documentclass[12pt]{amsart}
%\usepackage[margin=1in]{geometry}
\pagestyle{plain}

\usepackage{amsmath,amssymb}
\usepackage{enumerate}

\usepackage[shortlabels]{enumitem}
% Set spacing for all enumerate environments
\setlist[enumerate]{itemsep=6pt, topsep=6pt, parsep=0pt}

\usepackage{graphicx}

\usepackage[left=0.75in,right=1in,bottom=0.9in]{geometry}  
% \usepackage{mathptmx}      % use Times fonts if available on your TeX system
\usepackage[T1]{fontenc}

\usepackage[parfill]{parskip}
\renewcommand{\baselinestretch}{1}

%%%%%%%%%%%%%%%%%%%%%%%%%%%%%%%%%
% SECTION DIVIDERS
%%%%%%%%%%%%%%%%%%%%%%%%%%%%%%%%%
\newcommand{\divider}{
  \vspace{-1em} % Add space above
  \noindent
  \raisebox{-0.15ex}{\textbullet}%
  \hspace{0.5em}%
  \rule[0.5ex]{0.95\textwidth}{1.0pt}%
  \hspace{0.5em}%
  \raisebox{-0.15ex}{\textbullet}%
  \vspace{0em} % Adjust space below
}

\title{ESE 2030 QUIZZAM 1 : Fall 2025}
\author{prof-g}


\begin{document}
\maketitle

\begin{center}
    {\bf INSTRUCTIONS}
\end{center}

\begin{itemize}
\item 
No books, notes, or calculators. Use a writing utensil and logic. 

\item 
Please follow question instructions and fill in the bubble-sheet. 

\item 
Select one answer per problem.

\item 
Each problem has one best answer. In some problems, a wrong choice may be worth partial credit. 

\item 
Be careful to fill in the correct problem number. 

\item 
These questions are written carefully: if you detect an error, read closely and do your best.

\item 
If anything is unclear, use your best judgment. 

\item 
Cheating in any form will be dealt with severely. 
\end{itemize}

\newpage



% QUIZZAM 1: SELECTED PROBLEMS
% TOTAL QUESTIONS: 25
% COVERAGE: WEEKS 1-3

{\bf PROBLEM 1:}
% WEEK = 2
% TOPICS = Polynomial spaces, dimension, bases
Consider the vector space $\mathcal{P}_3$ of polynomials of degree at most 3. The dimension of this space is:
\begin{enumerate}[(A)]
\item 2
\item 3
\item 4
\item infinite
\item Cannot be determined without additional information
\end{enumerate}
% CORRECT ANSWER: (C): dim(P_3) = 4 (includes 1, x, x^2, x^3)
\divider

{\bf PROBLEM 2:}
% WEEK = 1
% TOPICS = Triangular systems, forward substitution, computational efficiency
For the lower triangular system $L\mathbf{x} = \mathbf{b}$ where $L = \begin{bmatrix} 2 & 0 & 0 \\ 1 & 3 & 0 \\ 4 & -1 & 5 \end{bmatrix}$, which statement about solving this system is most accurate?

\begin{enumerate}[(A)]
\item One should use row reduction to find the solution
\item The system may have no solutions, depending on $\mathbf{b}$
\item One should solve by forward substitution
\item The system may have infinitely many solutions depending on $\mathbf{b}$
\item One should solve by backward substitution
\end{enumerate}
% CORRECT ANSWER: (C): Lower triangular systems are solved by forward substitution - solve for $x_1$ from first equation, substitute into second to get $x_2$, then into third for $x_3$
\divider



{\bf PROBLEM 3:}
% WEEK = 3
% TOPICS = Linear transformations, kernel, matrix representation
Let $T: \mathbb{R}^3 \to \mathbb{R}^3$ be the linear transformation with matrix representation $A = \begin{bmatrix} 1 & 2 & 3 \\ 2 & 4 & 6 \\ -1 & -2 & -3 \end{bmatrix}$. What is the dimension of $\ker(T)$?
\begin{enumerate}[(A)]
\item 0 %(the transformation is injective)
\item 1 %(the nullity equals $3 - \text{rank}(A) = 3 - 2 = 1$)
\item 2 %(all rows are proportional, so rank is 1)
\item 3 %(the matrix maps everything to zero)
\item None of the above
\end{enumerate}
% CORRECT ANSWER: (C): All three rows are scalar multiples of $(1, 2, 3)$, so rank$(A) = 1$. By the Fundamental Theorem, $\dim(\ker(T)) = 3 - 1 = 2$
\divider

\newpage

{\bf PROBLEM 4:}
% WEEK = 2
% TOPICS = Vector spaces, examples and non-examples, vector space axioms
Consider the following sets with standard addition and scalar multiplication operations. Which one forms a vector space?

\begin{enumerate}[(A)]
\item The set of all $2 \times 2$ matrices with determinant equal to 1
\item The set of all continuous functions $f: [0,1] \to \mathbb{R}$ such that $f(0) = f(1) = 0$
\item The set of all vectors $(x, y, z) \in \mathbb{R}^3$ such that $x + y + z = 1$
\item The set of all $2 \times 2$ matrices $A$ such that $A^2 = A$ (idempotent matrices)
\item The set $\{(x, y) \in \mathbb{R}^2 : x \geq 0, y \geq 0\}$ (the first quadrant)
\end{enumerate}
% CORRECT ANSWER: (B): This is a subspace of all continuous functions. The zero function satisfies f(0) = f(1) = 0, and the set is closed under addition and scalar multiplication
% PARTIAL CREDIT: None
\divider



{\bf PROBLEM 5:}
% WEEK = 3
% TOPICS = Rank, nullity, Fundamental Theorem of Linear Algebra
Let ${\mathcal D}: \mathcal{P}_3 \to \mathcal{P}_2$ be the differentiation operator taking cubic to quadratic polynomials. 
The rank and nullity of this operator are:
\begin{enumerate}[(A)]
\item rank = 3, nullity = 1
\item rank = 2, nullity = 2
\item rank = 3, nullity = 0
\item rank = 2, nullity = 1
\item rank = 4, nullity = 0
\end{enumerate}
% CORRECT ANSWER: (A): Image of D is all of P2 (dim 3), so rank=3. Domain P3 has dim 4. By rank-nullity, 4 = 3 + nullity, so nullity=1.
\divider


{\bf PROBLEM 6:}
% WEEK = 1
% TOPICS = Elementary matrices, determinants, invertibility
Let $E$ be an elementary matrix used for row reduction. Which statement must be TRUE?

\begin{enumerate}[(A)]
\item $\det(E) = 1$
\item $\det(E) = \pm 1$
\item $\det(E) \neq 0$
\item $\det(E^{-1}) = -\det(E)$
\item $\det(EA) = \det(A)$ for any matrix $A$
\end{enumerate}
% CORRECT ANSWER: (C): Elementary matrices are always invertible, so det(E) ≠ 0
% PARTIAL CREDIT: None
% TRICK ANSWER: (A): Only true for row addition operations
% TRICK ANSWER: (B): True for swaps and row additions, but not scaling
% TRICK ANSWER: (D): Should be det(E^{-1}) = 1/det(E)
% TRICK ANSWER: (E): Only true if det(E) = 1
\divider



{\bf PROBLEM 7:}
% WEEK = 3
% TOPICS = Coimage, cokernel, quotient spaces, kernel, image
Let $T: V \to W$ be a linear transformation. Which statement correctly describes the coimage and cokernel of $T$?

\begin{enumerate}[(A)]
\item $\text{coim } T = \ker T$ and $\text{coker } T = \text{im } T$
\item $\text{coim } T = V/\text{im } T$ and $\text{coker } T = W/\ker T$
\item $\text{coim } T = V/\ker T$ and $\text{coker } T = W/\text{im } T$
\item $\text{coim } T = \text{im } T$ and $\text{coker } T = \ker T$
\item $\text{coim } T = W/\ker T$ and $\text{coker } T = V/\text{im } T$
\end{enumerate}

% CORRECT ANSWER: (C): The coimage is the quotient of the domain by the kernel (V/ker T), and the cokernel is the quotient of the codomain by the image (W/im T)
% PARTIAL CREDIT: None
% TRICK ANSWER: (A): Students might think coimage and cokernel are just alternate names for kernel and image
% TRICK ANSWER: (B): Reverses which space gets quotiented by what - a plausible misconception
% TRICK ANSWER: (D): Another confusion between the fundamental spaces and their quotients
% TRICK ANSWER: (E): Mixes up domain and codomain in the quotient constructions
\divider


{\bf PROBLEM 8:}
% TOPICS = Permutation matrices, properties, PLU decomposition
Consider the permutation matrix $P = \begin{bmatrix} 0 & 1 & 0 \\ 0 & 0 & 1 \\ 1 & 0 & 0 \end{bmatrix}$. Which of the following statements about $P$ is TRUE?
\begin{enumerate}[(A)]
\item $P^2 = I$ %(applying $P$ twice gives the identity)
\item $P^3 = I$ %(applying $P$ three times gives the identity)
\item $P^{-1} = -P$ %(the inverse is the negative)
\item $P^{-1} = P$ %(reverses the permutation)
\item None of the above
\end{enumerate}
% CORRECT ANSWER: (B): This is a cyclic permutation that cycles rows: 1→2→3→1. After three applications, each row returns to its original position, so $P^3 = I$
\divider


{\bf PROBLEM 9:}
% WEEK = 2
% TOPICS = Linear independence, dimension, bases
Consider the set of Euclidean vectors 
\[
    S = \left\{
    \begin{pmatrix}1\\2\\0\end{pmatrix}, 
    \begin{pmatrix}0\\1\\1\end{pmatrix}, 
    \begin{pmatrix}1\\3\\1\end{pmatrix}
    \right\}
\]
Which statement about $S$ is most correct?

\begin{enumerate}[(A)]
\item $S$ spans $\mathbb{R}^3$
\item $S$ is linearly independent but does not span $\mathbb{R}^3$
\item $S$ is linearly dependent and does not span $\mathbb{R}^3$
\item $S$ is linearly dependent and spans $\mathbb{R}^3$
\item $S$ forms a basis for $\mathbb{R}^3$
\end{enumerate}
% CORRECT ANSWER: (C): The set is linearly dependent since (1,3,1) = (1,2,0) + (0,1,1), and spans only a 2-dimensional subspace
% PARTIAL CREDIT: (E): Correctly identifies that only two vectors are independent
\divider


{\bf PROBLEM 10:}
% WEEK = 2
% TOPICS = Subspaces, intersection and union of subspaces
Let $W_1$ and $W_2$ be subspaces of a vector space $V$. Which of the following is always a subspace of $V$?

\begin{enumerate}[(A)]
\item $W_1 \cup W_2$ (the union of $W_1$ and $W_2$)
\item $W_1 \cap W_2$ (the intersection of $W_1$ and $W_2$)
\item The set of all vectors in $W_1$ but not in $W_2$
\item The set of all vectors in $V$ but not in $W_1+W_2$
\item The quotient spaces $V/W_1$ and $V/W_2$
\end{enumerate}
% CORRECT ANSWER: (B): Intersection of subspaces is always a subspace
\divider

\newpage

{\bf PROBLEM 11:}
% QUESTION 11
% WEEK = 2
% TOPICS = Subspaces, dimension, bases, polynomial spaces
Let $V$ be the vector space of all $3 \times 3$ symmetric matrices. Which of the following is the dimension of $V$?

\begin{enumerate}[(A)]
\item 2
\item 3
\item 4
\item 6
\item None of the above
\end{enumerate}
% CORRECT ANSWER: (D) -- 3 on diagonal and 3 off-diagonal
\divider


{\bf PROBLEM 12:}
% WEEK = 3
% TOPICS = Linear transformations, kernel, image, rank-nullity theorem
Let $T: V \to W$ be a linear transformation between finite-dimensional vector spaces. If $\dim(V) = 5$ and $\dim(W) = 3$, which statement about $T$ must be TRUE?

\begin{enumerate}[(A)]
\item $T$ must be injective
\item $T$ must be surjective
\item The dimension of $\ker(T)$ is at least 2
\item The dimension of $\text{im}(T)$ is 3
\item None of the above
\end{enumerate}
% CORRECT ANSWER: (C): By rank-nullity, dim(ker(T)) = 5 - dim(im(T)) ≥ 5 - 3 = 2
\divider



{\bf PROBLEM 13:}
% WEEK = 1
% TOPICS = LU decomposition, matrix factorization, nonsingularity
Consider a square matrix $A$ that has an LU decomposition $A = LU$. Which of the following statements is necessarily TRUE?

\begin{enumerate}[(A)]
\item $A$ is nonsingular if and only if all diagonal entries of $L$ are nonzero
\item $A$ is singular if and only if the last diagonal entry of $U$ equals zero
\item The existence of an LU decomposition guarantees that $A$ is nonsingular
\item If $A$ is nonsingular, then both $L$ and $U$ must have positive diagonal entries
\item $A$ is nonsingular if and only if all diagonal entries of $U$ are nonzero
\end{enumerate}
% CORRECT ANSWER: (E): Since L has 1's on its diagonal (by convention), L is always invertible. Matrix A = LU is nonsingular iff U is nonsingular, which happens iff all diagonal entries of U are nonzero (for triangular matrices, nonsingularity is equivalent to having all nonzero diagonal entries)
% SMALL PARTIAL CREDIT: (B,D): 
\divider



{\bf PROBLEM 14:}
% WEEK = 1
% TOPICS = LU decomposition, matrix factorization, system solvability
For a square matrix $A$ with LU decomposition $A = LU$, which statement is always true?

\begin{enumerate}[(A)]
\item The matrix $L$ is invertible
\item The matrix $U$ is invertible  
\item The system $A\mathbf{x} = \mathbf{b}$ has a unique solution for any $\mathbf{b}$
\item The reduced row eschelon form (RREF) of $A$ is the identity
\item The determinant of $A$ equals the determinant of $L$
\end{enumerate}
% CORRECT ANSWER: (A): L is lower triangular with 1's on diagonal, hence always invertible
\divider



{\bf PROBLEM 15:}
% QUESTION 15
% WEEK = 2
% TOPICS = Basis, dimension, spanning sets, linear dependence
Let $V$ be a 4-dimensional vector space. Suppose $S = \{\mathbf{v_1, v_2, v_3, v_4, v_5}\}$ is a set of vectors in $V$ where $\{\mathbf{v_1, v_2, v_3, v_4}\}$ forms a basis for $V$. Which of the following must be TRUE about $\mathbf{v}_5$?

\begin{enumerate}[(A)]
\item $\mathbf{v}_5 = \mathbf{0}$ (the zero vector)
\item $\mathbf{v}_5$ can be uniquely written as a linear combination of $\mathbf{v_1, v_2, v_3, v_4}$
\item $\mathbf{v}_5$ is linearly independent from any three vectors chosen from $\{\mathbf{v_1, v_2, v_3, v_4}\}$
\item The set $\{\mathbf{v_2, v_3, v_4, v_5}\}$ cannot be a basis for $V$
\item The set $S$ is linearly independent
\end{enumerate}
% CORRECT ANSWER: (B): Since {v_1, v_2, v_3, v_4} is a basis, it spans V. Therefore v_5 can be uniquely expressed as a linear combination of the basis vectors
% PARTIAL CREDIT: None
\divider



{\bf PROBLEM 16:}
% WEEK = 1
% TOPICS = Linear systems, row reduction, pivots
Consider the system $A\mathbf{x} = \mathbf{b}$ where $A$ is a $3 \times 3$ matrix. If row reduction of the augmented matrix $[\, A\, |\, \mathbf{b}\, ]$ yields exactly two pivots, which statement must be TRUE?

\begin{enumerate}[(A)]
\item The system has a unique solution
\item The system has infinitely many solutions if it has any solution at all
\item The system is inconsistent
\item The matrix $A$ is invertible
\item None of the above
\end{enumerate}
% CORRECT ANSWER: (B): With 2 pivots in a 3x3 system, there's one free variable, so if consistent, infinitely many solutions
\divider



{\bf PROBLEM 17:}
% WEEK = 3
% TOPICS = Linear transformations, additivity, homogeneity, definition
Which one of the following is NOT a linear transformation?

\begin{enumerate}[(A)]
\item $T: \mathbb{R}^2 \to \mathbb{R}^2$ defined by $T(x,y) = (x-y, 0)$
\item $S: \mathcal{P}_2 \to \mathcal{P}_2$ defined by $S(p(x)) = xp'(x)$ where $p'$ denotes the derivative
\item $R: \mathbb{R}^3 \to \mathbb{R}$ defined by $R(x,y,z) = 2x - 3y + z$
\item $F: \mathbb{R}^2 \to \mathbb{R}^3$ defined by $F(x,y) = (x+1, y, x-y)$
\item $G: \mathbb{R}^{2\times 2} \to \mathbb{R}$ defined by $G(A) = \text{trace}(A)$ (trace is the sum of diagonal terms)
\end{enumerate}
% CORRECT ANSWER: (D): F is not linear because F(0,0) = (1,0,0) ≠ (0,0,0), violating T(0) = 0. Also fails additivity: F(x₁+x₂, y₁+y₂) = (x₁+x₂+1, y₁+y₂, x₁+x₂-y₁-y₂) ≠ F(x₁,y₁) + F(x₂,y₂)
% PARTIAL CREDIT: None
% TRICK ANSWER: (A): Students might think the zero in second component makes it non-linear
% TRICK ANSWER: (B): Product of x and p'(x) might seem non-linear, but this IS linear in p
% TRICK ANSWER: (C): Students might incorrectly think maps to different dimension can't be linear
% TRICK ANSWER: (E): Students unfamiliar with trace might assume it's non-linear
\divider


{\bf PROBLEM 18:}
% WEEK = 3
% TOPICS = Differentiation operator, kernel, polynomial spaces
The Fundamental Theorem of Linear Algebra says that for a linear transformation $T:V\to W$ 

\begin{enumerate}[(A)]
\item The kernel and cokernel are isomorphic
\item The domain and codomain are isomorphic
\item The coimage and cokernel are isomorphic
\item The kernel and image are isomorphic
\item None of the above
\end{enumerate}
% CORRECT ANSWER: (E) none of the above -- it's the image and coimage that are isomorphic
\divider


{\bf PROBLEM 19:}
% WEEK = 3
% TOPICS = Isomorphism, dimension, vector space equivalence
Consider an $m$-by-$n$ matrix $A$ and the associated linear transformation $T_A:\mathbb{R}^n\to\mathbb{R}^m$. Which of the following is true? 

\begin{enumerate}[(A)]
\item $A$ is nonsingular if $T_A$ is injective
\item The column space of $A$ is the coimage of $T_A$
\item The null space of $A$ is the kernel of $T_A$
\item If $T_A$ is surjective then $A$ is nonsingular
\item If $T_A$ is injective then $A$ is singular 
\end{enumerate}
% CORRECT ANSWER: (C)
% PARTIAL CREDIT (B)
\divider


{\bf PROBLEM 20:}
% WEEK = 2  
% TOPICS = Linear independence, span, basis
Let $V$ be a vector space and $S = \{\mathbf{v_1, v_2, ..., v_n}\}$ be a set of vectors in $V$. If every vector in $V$ can be written as a linear combination of vectors in $S$, then which statement is true by definition?

\begin{enumerate}[(A)]
\item $S$ is a basis for $V$
\item $S$ is linearly independent
\item The dimension of $V$ is at most $n$ 
\item The dimension of $V$ is exactly $n$
\item $S$ spans $V$
\end{enumerate}
% CORRECT ANSWER: (E)
% PARTIAL CREDIT: (C) -- true, but not by definition
\divider


{\bf PROBLEM 21:}
% WEEK = 2
% TOPICS = Linear independence, polynomial spaces
In $\mathcal{P}_2$, consider the set $S = \{1+x, x+x^2, 1+x^2\}$. This set is:
\begin{enumerate}[(A)]
\item Linearly independent and spans $\mathcal{P}_2$
\item Linearly independent but does not span $\mathcal{P}_2$
\item Linearly dependent and spans $\mathcal{P}_2$
\item Linearly dependent and does not span $\mathcal{P}_2$
\item None of the above
\end{enumerate}
% CORRECT ANSWER: (A): The three polynomials are independent and dim(P_2)=3
\divider



{\bf PROBLEM 22:}
% WEEK = 3
% TOPICS = Coimage, quotient spaces, fundamental theorem
Let $T: \mathbb{R}^5 \to \mathbb{R}^6$ be a linear transformation with $\dim(\ker(T)) = 2$. The cokernel of $T$ is isomorphic to which space?

\begin{enumerate}[(A)]
\item $\mathbb{R}^1$ 
\item $\mathbb{R}^6$ 
\item $\mathbb{R}^4$ 
\item $\mathbb{R}^3$
\item $\ker(T)$
\end{enumerate}
% CORRECT ANSWER: (D): By the fundamental theorem of linear algebra
\divider


\newpage



{\bf PROBLEM 23:}
% TOPICS = Quotient spaces, equivalence classes, vector equivalence
Let $V = \mathbb{R}^3$ and let $W$ be the subspace defined by the $xy$-plane, i.e., $W = \{(x, y, 0) \mid x, y \in \mathbb{R}\}$. An equivalence relation is defined on $V$ where two vectors $\mathbf{u}$ and $\mathbf{v}$ are equivalent if their difference $\mathbf{u - v}$ is an element of $W$.

Given the vector $\mathbf{v} = (3, 5, 2)^T$, which of the following vectors belongs to the equivalence class $[\mathbf{v}]$?

\begin{enumerate}[(A)]
\item $(3, 5, 0)^T$
\item $(6, 10, 4)^T$
\item $(-1, -2, 2)^T$
\item $(0, 0, 0)^T$
\item None of the above
\end{enumerate}
% CORRECT ANSWER: (C): Two vectors $u, v$ are in the same equivalence class if $u-v \in W$. We check the difference for each option: $v - (-1, -2, 2) = (3 - (-1), 5 - (-2), 2 - 2) = (4, 7, 0)$. Since this result has a zero z-component, it lies in W.
% PARTIAL CREDIT: None
% TRICK ANSWER: (A): This is the projection of v onto W, a common point of confusion.
% TRICK ANSWER: (B): This is a scalar multiple of v; their difference $v - 2v = -v$, which is not in W.
\divider



{\bf PROBLEM 24:}
% WEEK = 1
% TOPICS = Elementary matrices, row operations, matrix multiplication order
If $B$ is obtained from matrix $A$ by first swapping rows 1 and 2, then adding 2 times row 1 to row 3, which equation correctly relates them?

\begin{enumerate}[(A)]
\item $B = E_2 E_1 A$ where $E_1$ swaps rows 1 and 2, $E_2$ adds 2 times row 1 to row 3
\item $B = E_1 E_2 A$ where $E_1$ swaps rows 1 and 2, $E_2$ adds 2 times row 1 to row 3
\item $B = A E_1 E_2$ where $E_1$ swaps rows 1 and 2, $E_2$ adds 2 times row 1 to row 3
\item $B = (E_2 E_1)^{-1} A$
\item None of the above
\end{enumerate}
% CORRECT ANSWER: (A): Row operations are applied left-to-right as E2(E1(A))
\divider



{\bf PROBLEM 25:}
% WEEK = 2
% TOPICS = Subspace verification, closure properties
Consider the subset $S = \{(x, y, z)^T \in \mathbb{R}^3 : x + y + z = 1\}$. Under standard vector addition and scalar multiplication, which property does $S$ violate to be a subspace?

\begin{enumerate}[(A)]
\item Closure under addition 
\item Closure under scalar multiplication 
\item Contains the zero vector 
\item Both (A) and (C)
\item All three: (A), (B), and (C)
\end{enumerate}
% CORRECT ANSWER: (E): The set $S$ is an affine space (plane not through origin). It fails all three: $(0,0,0) \notin S$; if $u, v \in S$ then $u+v$ has sum of coordinates = 2; if $u \in S$ and $c \neq 1$, then $cu$ has sum = $c$
% PARTIAL CREDIT: (D): Correctly identifies two violations but misses scalar multiplication
% TRICK ANSWER: (A): Only identifies one violation
% TRICK ANSWER: (B): Only identifies one violation
% TRICK ANSWER: (C): Only identifies the zero vector issue
\divider



\end{document}